\documentclass[11pt,a4paper]{article}
\usepackage[UTF8]{ctex}
\usepackage{amsmath,amssymb}
\usepackage{graphicx}
\usepackage{hyperref}
\usepackage{booktabs}
\usepackage{geometry}
\usepackage{longtable}
\usepackage{array}
\usepackage{enumitem}
\geometry{margin=2.2cm}
\hypersetup{colorlinks=true,linkcolor=blue,urlcolor=blue}
\graphicspath{{./}}
\title{\textbf{dev74\_1 —— 服务器运行指南（MG vs BiStabCG 加速比 > 1.5）}}
\author{ZhangXin \\ IMP, CAS}
\date{2026-08-13}
\begin{document}
\maketitle

\begin{abstract}
服务器（512G 内存 / 32G 显存，dev73\_5 实测 GPU：Tesla V100-SXM2-32GB）上运行
Clover MultiGrid 相对 BiStabCG 的加速比验证流程，目标 \texttt{speedup\_min $>=$ 1.5}。
依据：dev73\_5 在 V100-32G 实测 8x8x8x16 = \textbf{2.43x}。本指南以 Step 1 为
\textbf{强制闸门}：8x8x8x16 达标后才进入大格子步骤。一键执行：
\texttt{RUN=1 bash examples/qcu/mg\_dev74\_1\_server.sh}。
\end{abstract}

\section{环境准备（一次性）}
\begin{enumerate}
  \item 代码与构建：\texttt{git clone gitee.com/zhangxin8069/PyQCU}；\texttt{source ./env.sh}；
    \texttt{bash ./build.sh}（C++ CUDA 后端）；\texttt{bash ./install.sh}（Cython 扩展）。
  \item 自检：\texttt{nvidia-smi}（确认 32G 显存）；\texttt{torch.cuda.is\_available()}；
    \texttt{from pyqcu.cuda import qcu}（Cython 桥）。
  \item （可选）同步本地 \texttt{logs/nullvec\_cache} 避免重复构建粗算子。
\end{enumerate}

\section{Step 1 —— 快速验证（强制闸门，预期 2.4x）}
\begin{verbatim}
python examples/qcu/dev74/mg_dev74_clean.py --lattice 8 8 8 16 --prec c64 \
    --levels 2 --restart 10 --ct 1e5 --cmi 15 --pairs 5
python examples/qcu/dev74/mg_dev74_1_check.py --gate 1.5 --label "Step1 gate" \
    --file logs/dev74/dev74_clean_L8x8x8x16*.json
\end{verbatim}
闸门失败时：检查 GPU 占用/时钟节流（\texttt{nvidia-smi -q -d CLOCK}）、
\texttt{logs/clover\_multigrid.log} 异常、换卡重试。

\section{Step 2 —— 参数优化扫描（8x16x16x16）}
\texttt{python examples/qcu/mg\_dev74\_1\_sweep.py --lattice 8 16 16 16}
自动扫描 9 配置。参数优先级（本地与 V100 数据交叉一致）：

\begin{center}
\begin{tabular}{lcc}
\hline
参数 & 推荐 & 效果 \\
\hline
levels & \textbf{3L} & 最强优化项（V100: 1.32x vs 2L 1.16x） \\
restart r & \textbf{20} & 少进粗层（V100: 1.26x） \\
cmi & 15--200 & 影响小 \\
ct & 1e5 & 大容差减少粗层迭代 \\
\hline
\end{tabular}
\end{center}

\section{Step 3/4 —— 大格子（预算来自 dev74 实测校准模型）}
\begin{itemize}
  \item \textbf{16x32x32x32}：cold 28.4GB / warm 14.1GB，单卡全流程可行；
    粗算子构建首次 1--2 小时（\texttt{--build cpp} 快 2 倍），缓存命中后秒级。
  \item \textbf{16x32x32x64}：warm 28.2GB 可行；cold 56.9GB 超 32G——
    \textbf{分阶段}：先 \texttt{--build cpp --pairs 1} 构建缓存，再 warm 测量。
  \item \textbf{24x32x32x64}：warm 42.3GB 单卡不可行，需多卡 MPI（后续工作）。
\end{itemize}

\section{Step 5 —— 汇总断言与一键流程}
\texttt{python examples/qcu/mg\_dev74\_1\_check.py --gate 1.5 --file logs/dev74\_1\_sweep.json}
一键：\texttt{RUN=1 bash examples/qcu/mg\_dev74\_1\_server.sh}
（Step 0--5 全自动，任一步断言失败即停止，退出码非零）。

\section{故障排查速查}
\begin{center}
\begin{tabular}{ll}
\hline
现象 & 排查 \\
\hline
闸门失败 & GPU 占用/时钟节流/驱动；查 clover\_multigrid.log \\
16x32x32x64 cold OOM & 阶段 A/B 分离；--build cpp；确认空闲显存 $\ge$ 30G \\
构建太久 & nullvec 缓存复用；--build cpp \\
check 无数据 & 先跑 Step 1/2 生成 json \\
\hline
\end{tabular}
\end{center}

\section{参数约定}
mass=0.05, atol=$10^{-6}$, gauge\_seed=42, $\kappa=1/(2m+8)$, E=48,
NV\_ITERS=2, MG\_GRID=[2,2,2,2]，参考=applyCloverBistabCgQcu（VERBOSE=0）；
独立进程 + ref/mg 交叉计时 + min of N pairs。
\end{document}
